\chapter{Derivation of the plane-strain thermoelastic formulation}
\label{ch:derivations}

The relations used for design estimates and for the solid-solver sanity checks are derived here
from the constitutive law \eqref{eq:constitutive}. Where a simpler expression circulating in the
working notes during preparation was found to be in error, the corrected expression is derived and
the correction is indexed in the discrepancy register (Appendix~B, Table~\ref{tab:register}). The
underlying constitutive theory is classical \cite{boley1960theory,timoshenko1970elasticity}; the
specific reductions for the present geometry are derived here and are not attributed to the
literature.

\section{Plane-strain reduction of the constitutive law}
\label{sec:deriv:planestrain}
Write \eqref{eq:constitutive} componentwise with $e=\tr\beps=\eps_{xx}+\eps_{yy}$ (plane strain).
In compliance form, $\eps_{zz}=[\sigma_{zz}-\nu(\sigma_{xx}+\sigma_{yy})]/E+\alpha\theta=0$ gives
\eqref{eq:sigmazz} directly. In stiffness form, $\sigma_{zz}=\lambda e-\beta\theta$; substitution
of $\lambda=E\nu/[(1+\nu)(1-2\nu)]$, $\mu_L=E/[2(1+\nu)]$, and the in-plane stress--strain relations
confirms that the two expressions coincide identically (algebra in Appendix~\ref{app:arith}, cell
A2).

For a strip with traction-free faces ($\sigma_{yy}=0$ locally at each $x$), elimination of
$\eps_{yy}$ through $\sigma_{yy}=2\mu_L\eps_{yy}+\lambda e-\beta\theta=0$ yields
$e=(2\mu_L\eps_{xx}+\beta\theta)/(\lambda+2\mu_L)$; back-substitution into
$\sigma_{xx}=2\mu_L\eps_{xx}+\lambda e-\beta\theta$, using
$\lambda+2\mu_L=E(1-\nu)/[(1+\nu)(1-2\nu)]$, produces the one-dimensional wide-beam law
\begin{equation}
\sigma_{xx}=\frac{E}{1-\nu^2}\Big[\eps_{xx}-(1+\nu)\,\alpha\theta\Big].
\label{eq:widebeam}
\end{equation}
Form \eqref{eq:widebeam} makes effective constants $E/(1-\nu^2)$ and $(1+\nu)\alpha$ explicit; both
appear in the curvature and buckling relations below and distinguish the plane-strain strip from a
narrow beam \cite{boley1960theory}. Dimensions: $E\,[\mathrm{Pa}]$ multiplies a dimensionless
strain; $(1+\nu)\alpha\theta$ is dimensionless.

\section{Mean-stress identity}
\label{sec:deriv:meanstress}
From \eqref{eq:meanstress} and \eqref{eq:sigmazz},
$\sigma_m=(\sigma_{xx}+\sigma_{yy}+\sigma_{zz})/3$ takes the stress form
\begin{equation}
\sigma_m=\frac{1+\nu}{3}\big(\sigma_{xx}+\sigma_{yy}\big)-\frac{E\alpha}{3}\,\theta .
\tag{\ref{eq:meanstress}a}
\end{equation}
Summing the in-plane constitutive relations gives
$\sigma_{xx}+\sigma_{yy}=2(\lambda+\mu_L)e-2\beta\theta$; inserting this in (\ref{eq:meanstress}a)
and using $2(1+\nu)(\lambda+\mu_L)=E/(1-2\nu)=3K$ and $2(1+\nu)\beta+E\alpha=3\beta$
(Appendix~\ref{app:arith}) yields the strain form
\begin{equation}
\sigma_m=K\big(e-3\alpha\theta\big),
\tag{\ref{eq:meanstress}b}
\end{equation}
i.e.\ mean stress equals bulk modulus $\times$ mechanical volumetric strain. The two forms are
algebraically identical and serve as mutual checks in post-processing: the first requires complete
Cauchy components (never the mechanical part alone), the second requires only
$\tr\beps$ and $\theta$.

\section{Reference equilibrium states (sanity checks)}
\label{sec:deriv:sanity}
Three elementary states are used to test the constitutive implementation.
\textbf{(a) Fully constrained body, uniform rise $\theta$ ($\beps\equiv0$).} Then
$\bsig=-\beta\theta\bI$, so $\sigma_{xx}=\sigma_{yy}=\sigma_{zz}=-\beta\theta$ and
\begin{equation}
\sigma_m=-3K\alpha\theta=-\beta\theta .
\label{eq:vca1}
\end{equation}
\textbf{(b) Axially restrained plane-strain strip with traction-free faces, uniform $\theta$.}
Constraints $\eps_{xx}=0$, $\sigma_{yy}=0$, $\eps_{zz}=0$ give, from \eqref{eq:widebeam} and
\eqref{eq:sigmazz},
\begin{equation}
\sigma_{xx}=\sigma_{zz}=-\frac{E\alpha\theta}{1-\nu},\qquad
\eps_{yy}=\frac{1+\nu}{1-\nu}\,\alpha\theta,\qquad
\sigma_m=-\frac{2E\alpha\theta}{3(1-\nu)} .
\label{eq:vcb1}
\end{equation}
The mean stress of case (b), $-2E\alpha\theta/[3(1-\nu)]$, differs from $-3K\alpha\theta$: the two
values belong to different restraint conditions and were conflated in an earlier draft (Appendix~B,
register D-12); both are retained, correctly paired, as unit-test expectations in
Chapter~\ref{ch:validation}.
\textbf{(c) Free plane-strain strip, uniform $\theta$.} With
$\sigma_{xx}=\sigma_{yy}=0$, equation \eqref{eq:sigmazz} forces $\sigma_{zz}=-E\alpha\theta$, and
\begin{equation}
\sigma_m=-\frac{E\alpha\theta}{3}\neq0 .
\label{eq:freemean}
\end{equation}
Consistency with (\ref{eq:meanstress}b) follows from
$e=\eps_{xx}+\eps_{yy}=2(1+\nu)\alpha\theta$, which gives
$K\,[2(1+\nu)-3]\alpha\theta=E(2\nu-1)\alpha\theta/[3(1-2\nu)]=-E\alpha\theta/3$.
This value is the plane-strain analogue of the stress-free three-dimensional expansion and is the
correct expectation for the strip test (register D-2).

\section{Free thermal curvature of the strip}
\label{sec:deriv:curvature}
Let the through-thickness temperature be linear,
$\theta(y)=\bar\theta+\Gamma y$ measured from the mid-plane with $\Gamma=\Delta T_t/d$, and the
strip be free of external loads with traction-free faces. With
$\eps_{xx}=\eps_0+\kappa y$ in \eqref{eq:widebeam}, zero axial force resultant gives
$\eps_0=(1+\nu)\alpha\bar\theta$, and zero moment resultant,
$\int_{-d/2}^{d/2}\sigma_{xx}\,y\,\mathrm dy=0$, gives the free thermal eigen-curvature
\begin{equation}
\kappa_{\mathrm{th}}=(1+\nu)\,\alpha\,\frac{\Delta T_t}{d}
=(1+\nu)\,\alpha\,\frac{q^{''}}{k_s},
\label{eq:kappa}
\end{equation}
where the second form follows from one-dimensional steady conduction, $\Delta T_t=q^{''}d/k_s$
\cite{carslaw1947conduction}; \eqref{eq:kappa} is the plane-strain specialisation of the classical
thermal-curvature result \cite{boley1960theory}. At $\nu=0$ the narrow-beam value
$\kappa=\alpha\Delta T_t/d$ is recovered; a factor $(1+\nu)/(1-\nu)$ would pertain to biaxially
constrained curvature, not to the present free bending, and its appearance in an earlier draft is
recorded in register D-1. Two consequences matter for the design. First,
$\kappa_{\mathrm{th}}$ is independent of $d$: the through-wall gradient adjusts exactly to cancel
the thickness. Second, a clamped--clamped strip at uniform $\kappa_{\mathrm{th}}$ carries a
constant corrective moment $M_0=D_s\kappa_{\mathrm{th}}$ and does not bow at all
($w\equiv0$); pinned supports, or axial non-uniformity of $\kappa_{\mathrm{th}}(x)$, are required
for the bowing mechanism to exist, which fixes the support choice of Assumption~A9.

\section{Pinned--pinned thermal bowing}
\label{sec:deriv:deflection}
For the reduced wall equation \eqref{eq:beam} with no mechanical load, the curvature relation is
$w^{''}(x)=-\kappa_{\mathrm{th}}(x)$, with $w$ positive away from the channel
(Section~\ref{sec:signs}). For constant $\kappa_{\mathrm{th}}$ and $w(0)=w(L)=0$,
integration gives
\begin{equation}
w(x)=\frac{\kappa_{\mathrm{th}}}{2}\,x(L-x),\qquad
w_{\max}=\frac{\kappa_{\mathrm{th}}L^2}{8} .
\label{eq:wmax}
\end{equation}
Dimensions: $\si{\per\metre}\times\si{\square\metre}=\si{\metre}$. Relation \eqref{eq:wmax} is
exact within the thin-plate framework \emph{at negligible axial force}, which is the state of the
free-ended strip and of the zero-membrane benchmark specimen of Chapter~\ref{ch:validation}
($L/d=25$, keeping shear-deformation corrections below $1\%$
\cite{timoshenko1961stability}). For the axially constrained strip the thermal membrane force of
Section~\ref{sec:deriv:buckling} is not negligible at the baseline, and the bowing relation must be
replaced by the beam--column form of Section~\ref{sec:deriv:beamcolumn}.%

\section{Thermal membrane force and pinned--pinned buckling}
\label{sec:deriv:buckling}
With the axial displacement blocked at both ends (A9), the mean axial strain vanishes and
integration of \eqref{eq:widebeam} over the thickness yields the compressive membrane force per
unit depth
\begin{equation}
N_\theta=-\frac{Ed\,\alpha\bar\theta}{1-\nu} .
\label{eq:ntheta}
\end{equation}
A sign convention is fixed here and used throughout: $N_\theta$ is the \emph{signed} axial
resultant (tension positive, so thermal compression gives $N_\theta<0$), while the
beam--column parameter and every force ratio use the \emph{positive compressive magnitude}
$N_\theta^c=|N_\theta|$, so that $N_\theta^c/N_{\mathrm{cr}}\in(0,1)$ is sub-critical and
$\mu^2=N_\theta^c/D_s>0$. The compatibility statement of
Section~\ref{sec:deriv:beamcolumn} acts on the signed resultant; its correction term is
positive and therefore reduces the compression.
The pinned--pinned Euler threshold for the strip, using the wide-plate flexural rigidity of
\eqref{eq:beam}, is $N_{\mathrm{cr}}=\pi^2D_s/L^2$ \cite{timoshenko1961stability}; the wide-plate
value is the consistent one here \emph{because plane strain suppresses the antielastic
(out-of-plane) curvature mode}, so the strip cannot shed bending stiffness into the out-of-plane
direction --- a strip free to anticlastically buckle would carry the reduced rigidity
$(1-\nu^2)\,Ed^3/12$. Equating
$|N_\theta|=N_{\mathrm{cr}}$ gives the critical mean temperature rise
\begin{equation}
\bar T_{\mathrm{cr}}=\frac{\pi^2d^2}{12\,\alpha(1+\nu)L^2} .
\label{eq:tcr}
\end{equation}
The clamped--clamped eigenvalue $4\pi^2D_s/L^2$ would give a threshold four times higher; the
pinned value \eqref{eq:tcr} is the one consistent with the adopted supports and replaces the
clamped-based estimate used in an early design note (Appendix~B, register D-4). The admissibility
rule is $\bar\theta_{\max}\le\bar T_{\mathrm{cr}}/M$, with margin $M=2$ used in
Chapter~\ref{ch:baseline}. The full-edge pin idealization of Assumption~A9 stiffens the ends relative to ideal point pins by
$0.2\%$ at the baseline aspect ratio ($L/d=10$): the plane-strain Q8 computation of register D-26
gives $1.0000$ of $\kappa_{\mathrm{th}}L^2/8$ for point pins and $0.998$ for the full-edge pin,
against $-0.02$ to $-0.06$ for a clamp. The earlier reading of this end effect as ``a few per cent''
is superseded by that measurement, the check is complete, and no separate end-compliance correction
is applied to the catalogue.

\section{Beam--column bowing at finite membrane force}
\label{sec:deriv:beamcolumn}
The linear bowing of Section~\ref{sec:deriv:deflection} and the membrane force of
Section~\ref{sec:deriv:buckling} are not independent: a pinned strip bearing a compressive
force $N_\theta^c=|N_\theta|$ couples its transverse deflection to that force at leading order,
which is precisely the mechanism behind the threshold \eqref{eq:tcr}. The small-strain,
moderate-rotation (von K\'arm\'an) wall model for this strip consists of the transverse
equilibrium equation \emph{together with} an axial compatibility condition that links the
membrane force to the deflection. What is adopted here is the \emph{reduced (frozen-force)
beam--column equation}: the transverse member alone, retaining the full geometric-stiffness
contribution of the compressive membrane force but freezing that force at its blocked-state
value $N_\theta^c$:
\begin{equation}
D_s\,w^{''''}(x)+N_\theta^c\,w^{''}(x)=0,
\qquad
w(0)=w(L)=0,\qquad
w^{''}(0)=w^{''}(L)=-\kappa_{\mathrm{th}},
\label{eq:beamcol}
\end{equation}
where the end condition on $w^{''}$ states that the \emph{total} bending moment
$M=D_s\,(w^{''}+\kappa_{\mathrm{th}})$ vanishes at the pins. $M$ is taken positive when it produces
sagging, $M=-\int_{-d/2}^{d/2}\sigma_{xx}\,y\,\mathrm dy$, so the physical first moment of the stress
is $-M$; the vanishing condition itself is unchanged (register D-21). With $\mu^2=N_\theta^c/D_s$ the
exact solution is
\begin{equation}
w(x)=\frac{\kappa_{\mathrm{th}}}{\mu^2}
\left[\frac{1-\cos\mu L}{\sin\mu L}\,\sin\mu x+\cos\mu x-1\right],
\label{eq:beamcolsol}
\end{equation}
which recovers the parabola \eqref{eq:wmax} as $\mu L\to0$ and diverges at $\mu L=\pi$, i.e.\
exactly at the Euler threshold \eqref{eq:tcr}; the buckling limit and the deflection law are
therefore two limits of one relation, as they must be. The fully consistent von K\'arm\'an
formulation would additionally update the force through axial compatibility: with both end
edges blocked, $\langle u'\rangle=0$ in the mean strain
$\langle\eps_x\rangle=\langle u'\rangle+\tfrac12\langle w^{'}{}^2\rangle-(1+\nu)\alpha\bar\theta$,
so the signed resultant acquires the arc-stretching correction
\begin{equation}
N_\theta(w)\;=\;N_\theta\;+\;\frac{Ed}{2(1-\nu^2)}\,\langle w^{'}{}^2\rangle ,
\label{eq:nupdate}
\end{equation}
and the deflection problem becomes nonlinear. Equation \eqref{eq:beamcol} is the first iterate
of that system --- the force frozen at $N_\theta^c=|N_\theta|$ --- and its accuracy against the
consistent solution is quantified below rather than assumed. At the baseline parameters
($N_\theta^c/N_{\mathrm{cr}}=0.275$, $\mu L=1.65$):
\begin{equation}
w_{\max}=2.19\times10^{-4}\,\mathrm{m}=0.438\,H_0
\quad\text{(frozen force, adopted)},\qquad
\frac{w_{\max}^{\mathrm{lin}}}{H_0}=0.315
\quad\text{(linear, Eq.~\eqref{eq:wmax})},
\label{eq:beamcolbase}
\end{equation}
an amplification by the factor $1.39$, consistent at leading order with the beam--column
factor $1/(1-N_\theta^c/N_{\mathrm{cr}})=1.38$. The maximum slope is $w'_{\max}=0.027$~rad, the
bending fibre strain is $(\kappa_{\mathrm{geo}}-\kappa_{\mathrm{th}})d/2=9.9\times10^{-4}$ with
$\kappa_{\mathrm{geo}}=-w^{''}$. The stress-producing surface strain, $\varepsilon_{\mathrm{mech}}
=\varepsilon-\alpha_{\mathrm{eff}}\theta$ at $y=\pm d/2$, reaches $|\varepsilon_{\mathrm{mech}}|
=3.25\times10^{-3}$ on the fluid-side fibre (this is the fibre that carries $\sigma_{\max}$); the
kinematic measures are the blocked thermal strain $e_0=(1+\nu)\alpha\bar\theta=2.26\times10^{-3}$ and
the bending strain $\kappa_{\mathrm{geo}}d/2=3.09\times10^{-3}$, whose sum $e_0+\kappa_{\mathrm{geo}}
d/2=5.35\times10^{-3}$ is used below as a conservative upper bound on any surface strain. The
von K\'arm\'an hypotheses (strains $\ll1$, rotations $\ll1$~rad)
are comfortably satisfied, so \eqref{eq:beamcol} is an appropriate model level --- not a
large-deformation extension, but a reduction of the von K\'arm\'an system whose error must be
stated. Freezing the force is an approximation, not an identity: the deflection relaxes the
blocked-end mechanical strain $e_0=(1+\nu)\alpha\bar\theta=2.26\times10^{-3}$ by the
arc-stretching term $\tfrac12\langle w^{'}{}^2\rangle=1.39\times10^{-4}$ (circular-arc estimate
$1.6\times10^{-4}$), i.e.\ a first-iterate force relaxation of $6.1\%$ (register D-23). Solving the
compatibility relation \eqref{eq:nupdate} self-consistently (a scalar fixed point, since the
solution family \eqref{eq:beamcolsol} holds for every fixed force) gives a relaxation of
$5.9\%$ and lowers the baseline deflection to
$w_{\max}=2.14\times10^{-4}\,\mathrm{m}=0.428\,H_0$: the frozen-force value
\eqref{eq:beamcolbase} \emph{over}predicts the consistent deflection by $2.3\%$, a one-sided
(conservative) error that never changes sign. Within the admissible envelope the error remains
below $\approx4\%$ (the deformation and stress ceilings of Chapter~\ref{ch:baseline} bind at
$N_\theta^c/N_{\mathrm{cr}}=0.30$ and $0.33$), the mid-span moment and bending stress shift by
under $8\%$ in the same safe direction, and the buckling margin is exact --- the relaxation
term vanishes as the deflection vanishes, so the onset threshold is governed by the blocked
force and $\bar T_{\mathrm{cr}}$ is unaffected. The error grows rapidly only beyond
$N_\theta^c/N_{\mathrm{cr}}\approx0.5$ (it reaches $\approx17\%$ at the $M_b=2$ ceiling), which
is precisely where the escalation criterion of Chapter~\ref{ch:numerics} moves the computation
off the frozen-force form.

Three consequences fix the model specification. First, the solid model of the coupled problem is
specified at the reduced (frozen-force) von K\'arm\'an level --- the geometric-stiffness
contribution of the prescribed membrane state, of which the linear model is the small-force
limit (at $N_\theta^c/N_{\mathrm{cr}}\lesssim0.1$ the two differ by $\lesssim10\%$) --- while
the implementation re-evaluates the membrane state across coupled iterations
(Chapter~\ref{ch:numerics}), so the frozen-force form applies within each solve and the
consistency update \eqref{eq:nupdate} is recovered at the fixed point. Second, the earlier
design relation \eqref{eq:wmax} is retained for the free and zero-membrane states only, and the
baseline bowing is quoted from \eqref{eq:beamcolbase}. Third, the bending stress follows from
the beam--column moment: at mid-span $M=-6.4$~N\,m per unit depth, giving
$\sigma_{\mathrm{bend}}=6|M|/d^2=4.2$~MPa (the unreleased-moment scale $E\kappa_{\mathrm{th}}d/2
/(1-\nu^2)=9.0$~MPa quoted previously belongs to a fully restrained-curvature state and does not
occur in the pinned strip; Appendix~B, register D-17).

\section{Hydraulic loading and the deformation ratio}
\label{sec:deriv:pressure}
\textbf{Pressure drop.} Integrating the plane-Poiseuille profile
$u(y)=6U(y/H_0)(1-y/H_0)$ \cite{bird1960transport,shah1978laminar} through the momentum balance,
$\mathrm dp/\mathrm dx=\mu u^{''}=-12\mu U/H_0^2$, gives the slit pressure drop
\begin{equation}
\Delta p=\frac{12\mu UL}{H_0^2}
\qquad
\Big[\frac{(\mathrm{Pa\,s})(\mathrm{m/s})(\mathrm m)}{\mathrm{m^2}}=\mathrm{Pa}\Big] .
\label{eq:dp}
\end{equation}
A predecessor note carried the same expression with $H_0^3$ in the denominator, which is
dimensionally not a pressure (Pa/m) and produced an estimate three orders of magnitude too large;
the corrected value is computed in Chapter~\ref{ch:baseline} (register D-5).

\textbf{Pressure deflection.} For a \emph{pinned--pinned} strip under uniform transverse load
$\bar p$ acting in the same sense as the thermal bowing, the standard bending solution
\cite{timoshenko1961stability,timoshenko1970elasticity} gives
\begin{equation}
w_p=\frac{5\,\bar p L^4}{384D_s}.
\label{eq:wpdef}
\end{equation}
(The coefficient $5/384$ belongs to pinned supports; the value $1/384$ used in an earlier estimate
is the \emph{clamped--clamped} coefficient. See register D-13.) The load level requires a datum
statement, because the incompressible pressure \emph{level} on the wall is fixed by the outlet
datum of condition~B2 and by the casing environment of Assumption~A9, not by $\Delta p$ alone:
with the datum $p=0$ at the outlet and the casing at the same reference pressure, the gauge
pressure on the strip rises linearly from $p(L)=48$~Pa to $p(0)=120$~Pa (gradient
$12\mu U/H_0^2$ over the strip and the exit length), so the strip carries a linearly varying load
with mean
\begin{equation}
\bar p=\tfrac{1}{2}\,[p(0)+p(L)] = 84~\mathrm{Pa}
\label{eq:pbar}
\end{equation}
rather than the value $\Delta p/2=36$~Pa quoted in earlier drafts (register D-16); the
uniform-load formula \eqref{eq:wpdef} is then read as a design estimate for this mildly varying
load. With $\bar p$ of \eqref{eq:pbar}, $w_p=9.2\times10^{-8}$~m, and with $w_{\max}$ of
\eqref{eq:beamcolbase}, the deformation ratio
\begin{equation}
\Gamma_p=\frac{w_{\max}}{w_p}
=\frac{48\,\kappa_{\mathrm{th}}D_s}{5\,\bar p L^2}\,\frac{w_{\max}}{w_{\max}^{\mathrm{lin}}}
=2.4\times10^{3}\quad\text{at the baseline},
\label{eq:gammap}
\end{equation}
where $w_{\max}^{\mathrm{lin}}=\kappa_{\mathrm{th}}L^2/8$ and
$w_{\max}/w_{\max}^{\mathrm{lin}}=1.39$ is the beam--column amplification of
\eqref{eq:beamcolbase}. The linear-limit ratio is
$48\kappa_{\mathrm{th}}D_s/(5\bar pL^2)=1.7\times10^{3}$, which the amplification converts to the
quoted value (register D-22). The ratio measures the thermal against the hydraulic contribution to a
co-directional deflection. The numerator carries the beam--column amplification of the strip's own
membrane force while the denominator \eqref{eq:wpdef} does not, so $\Gamma_p$ of
\eqref{eq:gammap} is a \emph{mixed} ratio. It is retained because it answers the design question
asked --- how large is the thermal bow compared with the hydraulic deflection of the unstressed
strip --- but it is not a ratio of two responses of the same structure. The like-for-like
alternative, in which the same compressive force amplifies the pressure response, is
$w_p(N_\theta^c)=1.27\times10^{-7}$~m and gives $w_{\max}/w_p(N_\theta^c)=1.73\times10^{3}$,
essentially the linear-limit value, because the two deflections then share a common amplification.
Both ratios are $\gg5$, so the conclusion (thermal dominance) is unchanged, and it is insensitive
to the datum choice, since any plausible casing pressure shifts $\bar p$ by less than an order of
magnitude (register D-29).
Because $\kappa_{\mathrm{th}}$ is independent of $d$ while $D_s\propto d^3$ and
$\bar p\propto UH_0^{-2}$, the ratio grows with thicker walls and decreases with speed and with
reduced height: these are the control directions examined at the design stage.

\section{Time scales and the steady-state reduction}
\label{sec:steadytimes}
Four time scales govern the start-up problem (inputs of Table~\ref{tab:props}; arithmetic in
Appendix~\ref{app:arith}): the structural wave crossing
$t_{\mathrm{wave}}=d\sqrt{\rho_s/E}=1.8\times10^{-6}$ s; the fluid advection time
$t_{\mathrm{adv}}=L/U=0.6$ s; the fluid transverse diffusion time
$t_{f}=H_0^2/\kappa_f\approx1.8$ s; and the solid conduction time
$t_{s}=d^2/\kappa_s\approx51$ s. The ordering
$t_{\mathrm{wave}}\ll\{t_{\mathrm{adv}},t_{f}\}\lesssim t_s$ has three consequences. First,
mechanical equilibrium is quasi-static on all observable transport times, justifying
\eqref{eq:equilibrium}. Second, the steady coupled solution targeted by the baseline is the
$t\to\infty$ asymptote of the start-up dynamics, which is rate-limited by the solid conduction
time, not by the fluid. Third, in a future start-up computation the fluid transients may be
neglected at leading order ($t_f/t_s\approx0.04$) provided an a-posteriori check is made; the
thermoelastic deformation-heating source of Section~\ref{ch:equations} vanishes identically in the
steady problem, and its transient neglect is controlled by $\xi\approx2\times10^{-2}$
\cite{boley1960theory}.
